\documentclass[12pt,a4paper,twoside]{book}
% Language setup for Greek and English
\usepackage[english,greek]{babel}
\usepackage[utf8]{inputenc}
\usepackage[T1,LGR]{fontenc}
% Font setup
\usepackage{fontspec}
\setmainfont{Linux Libertine O}[Scale=0.9]
\setsansfont{Linux Biolinum}[Scale=0.9]
\setmonofont{DejaVu Sans Mono}[Scale=0.9]
% Other necessary packages
\usepackage{enumitem}
\usepackage{graphicx}
\usepackage{amsmath}
\usepackage{makeidx}
\usepackage{multicol}
\usepackage{multirow}
\usepackage{adjustbox}
\usepackage{amssymb}
\usepackage{stackengine}
\usepackage{ifthen}
\usepackage{array}
\usepackage{booktabs}
\usepackage{tcolorbox}
\tcbuselibrary{skins,breakable}
\usepackage{listings}
\usepackage{parskip}
\usepackage{float}
\usepackage{subcaption}
\usepackage{xcolor}
\usepackage{cancel}
\usepackage{pdfpages}
\usepackage{titlesec}

% Page & text layout
\usepackage{geometry}
\geometry{
   a4paper,
   top=2.5cm,
   bottom=2.5cm,
   left=2.5cm,
   right=2.5cm
}

% Headers & footers
\usepackage{fancyhdr}
\pagestyle{fancyplain}
\fancyhf{}

\fancypagestyle{plain}{%
    \fancyhf{}
    \renewcommand{\headrulewidth}{0pt}
    \renewcommand{\footrulewidth}{0pt}
}

\fancypagestyle{tocstyle}{%
    \fancyhf{}
    \fancyhead[LE]{\thepage}
    \fancyhead[CE]{\leftmark}
    \fancyhead[CO]{\rightmark}
    \fancyhead[RO]{\thepage}
    \renewcommand{\headrulewidth}{0pt}
}

\renewcommand{\sectionmark}[1]{\markboth{#1}{}}
\renewcommand{\subsectionmark}[1]{\markright{#1}}

\let\oldtableofcontents\tableofcontents
\renewcommand{\tableofcontents}{%
    \clearpage
    \thispagestyle{empty}
    \pagestyle{tocstyle}
    \oldtableofcontents
    \clearpage
    \pagestyle{fancy}
    \fancyhead[LE]{\thepage}
    \fancyhead[CE]{\leftmark}
    \fancyhead[RE]{ΚΕΦ. \thechapter}

    \fancyhead[LO]{\thesection}
    \fancyhead[CO]{\rightmark}
    \fancyhead[RO]{\thepage}

    \renewcommand{\headrulewidth}{0.4pt}

    \renewcommand{\chaptermark}[1]{\markboth{##1}{}}
    \renewcommand{\sectionmark}[1]{%
        \markright{##1}%
    }
    \renewcommand{\subsectionmark}[1]{}
}

\makeatletter
\renewcommand\sectionmark[1]{%
  \markright{%
    \ifnum \c@secnumdepth >\z@
      \thesection\quad
    \fi
    #1}}
\makeatother

\addtocontents{toc}{\protect\thispagestyle{plain}}

% Hyperlinks
\usepackage{hyperref}
\usepackage{bookmark}
\hypersetup{
    colorlinks=true,
    linkcolor=blue,
    citecolor=blue,
    filecolor=blue,
    urlcolor=customCodeBlue,
    unicode,
    pdftitle={Ανίχνευση Δικτυακών Εισβολών με Μηχανική Μάθηση},
    pdfsubject={Αναφορά Εργασίας Εξόρυξης Δεδομένων},
    pdfauthor={Βασίλειος Μπίτζας},
    pdfborderstyle={/S/U/W 1},
}

\definecolor{customRed}{HTML}{FF5F5A}
\definecolor{customYellow}{HTML}{FFBE2E}
\definecolor{customGreen}{HTML}{2ACA44}
\definecolor{customPurple}{HTML}{C477DB}
\definecolor{customBlue}{HTML}{052538}
\definecolor{customLightBlue}{HTML}{60ADEC}
\definecolor{customGray}{HTML}{A0A0A0}
\definecolor{customBlack}{HTML}{000000}

\definecolor{customBgGray}{HTML}{E0E0E0}
\definecolor{customCodeRed}{HTML}{C92A2A}
\definecolor{customCodeYellow}{HTML}{E67700}
\definecolor{customCodeGreen}{HTML}{087F23}
\definecolor{customCodePurple}{HTML}{7B1FA2}
\definecolor{customCodeBlue}{HTML}{1976D2}
\definecolor{customCodeBlack}{HTML}{1A1A1A}

\lstdefinestyle{myPythonStyle}{
    language=Python,
    backgroundcolor=\color{customBgGray},
    basicstyle=\color{customCodeBlack}\ttfamily\footnotesize,
    keywordstyle=\color{customCodePurple}\bfseries,
    keywordstyle=[2]\color{customCodeYellow}\bfseries,
    commentstyle=\color{customCodeGreen}\itshape,
    stringstyle=\color{customCodeGreen},
    identifierstyle=\color{customCodeBlue},
    morekeywords={self, None, True, False, format, abs,
                as, pass, return, if, elif, else, for,
                range, while, try, except, with, lambda,
                yield, global, nonlocal, assert, del,
                raise, in, is, and, or, not},
    morekeywords=[2]{os, pd, np, plt, sns, sklearn, pandas, numpy,
                    matplotlib, seaborn, RandomForestClassifier,
                    LogisticRegression, DecisionTreeClassifier,
                    KMeans, DBSCAN, StandardScaler, RobustScaler,
                    GridSearchCV, train_test_split, classification_report,
                    fit, transform, fit_transform, predict, score,
                    read_csv, read_parquet, to_parquet, drop, dropna,
                    drop_duplicates, corr, describe, value_counts,
                    select_dtypes, isna, isnull, info, head, shape},
    showstringspaces=false,
    breaklines=true,
    breakatwhitespace=false,
    postbreak={\space},
    breakautoindent=false,
    breakindent=0pt,
    resetmargins=true,
    keepspaces=true,
    showspaces=false,
    columns=flexible,
    numbers=left,
    numberstyle=\small\color{customBgGray},
    stepnumber=1,
    numbersep=8pt,
    xleftmargin=1.5em,
    tabsize=4,
    captionpos=t,
    firstnumber=auto,
    aboveskip=0em,
    belowskip=0em,
    frame=none,
    framerule=0pt,
    rulecolor=\color{customBgGray},
    framesep=2pt,
    escapeinside={\(*@}{@*\)},
    morecomment=[l]{\#},
    morestring=[b]',
    morestring=[b]",
    literate={
        {``}{{\textquotedbleft}}2
        {''}{{\textquotedbright}}2
        {`}{{\textquoteleft}}1
        {'}{{\textquoteright}}1
        {_}{{\_}}1
    },
}

\lstdefinestyle{myBashStyle}{
    language=bash,
    backgroundcolor=\color{customBgGray},
    basicstyle=\color{customCodeBlack}\ttfamily\footnotesize,
    keywordstyle=\color{customCodePurple}\bfseries,
    keywordstyle=[2]\color{customCodeYellow}\bfseries,
    commentstyle=\color{customCodeGreen}\itshape,
    stringstyle=\color{customCodeGreen},
    identifierstyle=\color{customCodeBlue},
    alsoletter={0123456789}
    morekeywords={if, then, else, elif, fi, for, while, do, done, in, case, esac,
                  function, select, until, break, continue, return, exit, shift,
                  declare, local, readonly, export, set, unset},
    morekeywords=[2]{echo, cat, grep, sed, cut, sort, uniq, wc,
                     mkdir, rm, cp, mv, ls, cd, pwd, chmod, find,
                     source, pip, python, jupyter, git},
    showstringspaces=false,
    frame=none,
    rulecolor=\color{customBgGray},
    breaklines=true,
    breakatwhitespace=false,
    postbreak={\space},
    breakautoindent=false,
    breakindent=0pt,
    resetmargins=true,
    keepspaces=true,
    showspaces=false,
    columns=flexible,
    numbers=left,
    numberstyle=\small\color{customBgGray},
    stepnumber=1,
    numbersep=8pt,
    xleftmargin=1.5em,
    tabsize=4,
    captionpos=t,
    firstnumber=auto,
    aboveskip=0em,
    belowskip=0em,
    framesep=2pt,
    escapeinside={\(*@}{@*\)},
    morecomment=[l]{\#},
    morestring=[b]",
    morestring=[b]',
    literate={
        {``}{{\textquotedbleft}}2
        {''}{{\textquotedbright}}2
        {`}{{\textquoteleft}}1
        {'}{{\textquoteright}}1
        {_}{{\_}}1
    },
}

\tcbset{
    myCustomStyle/.style={
        enhanced,
        breakable,
        colback=customBlue,
        colframe=customBgGray,
        listing only,
        fonttitle=\bfseries,
        interior style={
            top color=black!75,
            bottom color=black!75,
        },
        colbacktitle=black!75,
        coltitle=white,
        rounded corners,
        boxrule=1mm,
        drop shadow southeast,
        lefttitle=0pt,
        title={\hspace*{-1em}\textcolor{customRed}{● }\textcolor{customYellow}{● }\textcolor{customGreen}{●}\quad#1},
        attach title to upper={\vspace{2.3mm}},
        beforeafter skip=8pt,
        breaklines=true,
        breakatwhitespace=false,
    }
}

% Custom commands for language switching
\newcommand{\en}[1]{\foreignlanguage{english}{#1}}
\newcommand{\gr}[1]{\foreignlanguage{greek}{#1}}
\newcommand{\blue}[1]{\textcolor{blue}{#1}}

% Index
\makeindex
\setlength{\headheight}{14.5pt}

% Chapter/section spacing
\titleformat{\chapter}[display]
{\normalfont\huge\bfseries}{\chaptertitlename\ \thechapter}{15pt}{\Huge}
\titlespacing*{\chapter}{0pt}{-10pt}{20pt}

\titlespacing*{\section}{0pt}{10pt}{10pt}
\setlength{\parskip}{0.5em}

% Figure placement
\setcounter{topnumber}{2}
\setcounter{bottomnumber}{2}
\setcounter{totalnumber}{4}
\renewcommand{\topfraction}{0.85}
\renewcommand{\bottomfraction}{0.85}
\renewcommand{\textfraction}{0.15}
\renewcommand{\floatpagefraction}{0.7}

% Graphics path: figures live under outputs/figures/ at the repo root
\graphicspath{{../outputs/figures/}}

\begin{document}
\pagenumbering{arabic}
\let\cleardoublepage\clearpage

\selectlanguage{greek}

\title{Ανίχνευση Δικτυακών Εισβολών με Μηχανική Μάθηση}
\author{Βασίλειος Μπίτζας}
\date{\today}

%=============================================================================
% TITLE PAGE
%=============================================================================
\begin{titlepage}
    \vspace*{4cm}
    \begin{center}
        % \includegraphics[width=0.3\textwidth]{media/logo.png}\\
        % \vspace*{1cm}
        {\LARGE \textbf{\textit{Ανίχνευση Δικτυακών Εισβολών}}}\\
        \vspace*{0.3cm}
        {\LARGE \textbf{\textit{με Μηχανική Μάθηση}}}\\
        \vspace*{0.5cm}
        {\large Ανάλυση του Συνόλου Δεδομένων CIC-IDS-2017}\\
        \vspace*{0.3cm}
        {\large Εξόρυξη Δεδομένων \& Μηχανική Μάθηση}
        \vspace*{5cm}
        \\
        \vspace*{1cm}
        \begin{tabular}{l l l}
            \large Βασίλειος Μπίτζας & \large 1083796 & \large up1083796@ac.upatras.gr
        \end{tabular}

        \vspace*{1cm}

        {\normalsize Τμήμα Μηχανικών Ηλεκτρονικών Υπολογιστών \& Πληροφορικής}
        \\
        {\normalsize Πανεπιστήμιο Πατρών}
        \\
        \vspace*{0.5cm}
        {\normalsize \today}
    \end{center}
    \thispagestyle{empty}
\end{titlepage}

\tableofcontents
\listoffigures
\listoftables

%=============================================================================
% ΚΕΦΑΛΑΙΟ 1: ΕΙΣΑΓΩΓΗ
%=============================================================================
\chapter{Εισαγωγή}

\section{Πλαίσιο και Κίνητρο}

Η ραγδαία αύξηση της εξάρτησης των υποδομών από το διαδίκτυο έχει καταστήσει
την \textbf{ανίχνευση δικτυακών εισβολών} (Network Intrusion Detection) έναν
από τους πιο κρίσιμους τομείς της κυβερνοασφάλειας. Τα σύγχρονα συστήματα
ανίχνευσης εισβολών (Intrusion Detection Systems, IDS) δεν μπορούν πλέον να
βασίζονται αποκλειστικά σε \textit{signature-based} κανόνες, καθώς οι
επιτιθέμενοι προσαρμόζουν συνεχώς τις τεχνικές τους. Αντίθετα, οι
προσεγγίσεις που βασίζονται σε \textbf{μηχανική μάθηση} (machine learning)
μαθαίνουν μοτίβα από ιστορικά δεδομένα δικτυακής ροής και γενικεύουν σε νέες,
άγνωστες επιθέσεις.

Η παρούσα εργασία αναπτύσσει ένα ολοκληρωμένο pipeline εξόρυξης δεδομένων
πάνω στο \textbf{CIC-IDS-2017}, ένα από τα πιο ευρέως χρησιμοποιούμενα
δημόσια σύνολα δεδομένων για έρευνα IDS. Το σύνολο περιέχει \textbf{2.83
εκατομμύρια} εγγραφές δικτυακής ροής (flow records) κατανεμημένες σε μία
κανονική κλάση (\texttt{BENIGN}) και 14 διαφορετικούς τύπους επιθέσεων, από
DDoS και port-scanning μέχρι web attacks και εξαιρετικά σπάνια exploits
όπως το \texttt{Heartbleed}.

\section{Στόχοι Εργασίας}

Η εργασία δομείται γύρω από τρία συνδεδεμένα ερωτήματα:

\begin{enumerate}
    \item \textbf{Ερώτημα 1 - Exploratory Data Analysis.} Φόρτωση, καθαρισμός
    και στατιστική περιγραφή του CIC-IDS-2017. Αντιμετώπιση Inf/NaN, διπλοτύπων,
    ανισορροπίας κλάσεων και πολυσυγγραμμικότητας features. Εξαγωγή ενός
    καθαρού, cached snapshot για επαναχρησιμοποίηση στα Q2 και Q3.

    \item \textbf{Ερώτημα 2 - Ταξινόμηση.} Εκπαίδευση και σύγκριση τριών
    αλγορίθμων (Logistic Regression, Decision Tree, Random Forest) με δύο
    διαφορετικές ρυθμίσεις Grid Search (A και B) για hyperparameter tuning.
    Αξιολόγηση με macro-F1 και per-class μετρικές λόγω έντονης ανισορροπίας.

    \item \textbf{Ερώτημα 3 - Συσταδοποίηση.} Εφαρμογή K-Means, Hierarchical
    Clustering και DBSCAN σε unsupervised mode. Σύγκριση των συστάδων με τις
    γνωστές ετικέτες (external validation) για να εκτιμηθεί κατά πόσο η δομή
    των δεδομένων υποστηρίζει ανίχνευση επιθέσεων χωρίς labels.
\end{enumerate}

Η αναφορά αυτή επικεντρώνεται κυρίως στο \textbf{Ερώτημα 1}, το οποίο
αποτελεί το υπόβαθρο για τα υπόλοιπα δύο. Τα Q2 και Q3 θα προστεθούν ως
ξεχωριστά κεφάλαια καθώς ολοκληρώνονται.

\section{Περιγραφή Δεδομένων}

\subsection{CIC-IDS-2017 Dataset}

Το CIC-IDS-2017 δημιουργήθηκε από το Canadian Institute for Cybersecurity
(CIC) και κυκλοφορεί ως σύνολο από οκτώ αρχεία CSV, ένα για κάθε ημέρα και
χρονική ζώνη της εβδομαδιαίας καταγραφής:

\begin{itemize}
    \item \texttt{Monday-WorkingHours.pcap\_ISCX.csv} - μόνο κανονική κίνηση.
    \item \texttt{Tuesday-WorkingHours.pcap\_ISCX.csv} - Brute Force (FTP/SSH).
    \item \texttt{Wednesday-workingHours.pcap\_ISCX.csv} - DoS (Hulk, GoldenEye,
    slowloris, Slowhttptest, Heartbleed).
    \item \texttt{Thursday-WorkingHours-Morning-WebAttacks.pcap\_ISCX.csv} -
    Brute Force, XSS, SQL Injection.
    \item \texttt{Thursday-WorkingHours-Afternoon-Infilteration.pcap\_ISCX.csv} -
    Infiltration.
    \item \texttt{Friday-WorkingHours-Morning.pcap\_ISCX.csv} - Botnet.
    \item \texttt{Friday-WorkingHours-Afternoon-PortScan.pcap\_ISCX.csv} -
    Port Scan.
    \item \texttt{Friday-WorkingHours-Afternoon-DDos.pcap\_ISCX.csv} - DDoS.
\end{itemize}

Κάθε γραμμή αντιστοιχεί σε μία \textit{δικτυακή ροή} (flow), δηλαδή μια ακολουθία
πακέτων μεταξύ ενός ζεύγους IP/ports για ένα συγκεκριμένο πρωτόκολλο.
Τα features έχουν εξαχθεί με το εργαλείο \textbf{CICFlowMeter} και
περιλαμβάνουν:

\begin{itemize}
    \item \textbf{Βασικά} χαρακτηριστικά: \texttt{Destination Port},
    \texttt{Flow Duration}, \texttt{Total Fwd/Bwd Packets}.
    \item \textbf{Στατιστικά μήκους πακέτων}: mean, std, min, max σε forward
    και backward κατευθύνσεις.
    \item \textbf{Στατιστικά χρονικών διαστημάτων} (Inter-Arrival Times, IAT):
    \texttt{Flow IAT Mean/Std/Max/Min}, \texttt{Fwd IAT}, \texttt{Bwd IAT}.
    \item \textbf{Flags TCP}: \texttt{FIN/SYN/RST/PSH/ACK/URG Flag Count}.
    \item \textbf{Bulk και subflow} στατιστικά, χαρακτηριστικά
    active/idle χρόνου.
\end{itemize}

Η δεύτερη στήλη που αφορά την εργασία είναι το \textbf{Label} - η ετικέτα
της κλάσης (\texttt{BENIGN} ή όνομα επίθεσης).

\subsection{Συνοπτικά Μεγέθη}

\begin{table}[H]
    \centering
    \begin{tabular}{lc}
        \toprule
        \textbf{Μέγεθος}                 & \textbf{Τιμή}           \\
        \midrule
        Πηγαία αρχεία CSV                & 8                       \\
        Αρχικές γραμμές                  & 2.830.743               \\
        Αρχικές στήλες                   & 79 (+ \texttt{Label})   \\
        Κλάσεις (συμπ. \texttt{BENIGN})  & 15                      \\
        Μετά από καθαρισμό (γραμμές)     & 2.572.640               \\
        Μετά από feature selection       & 47 features + Label     \\
        Cached Parquet snapshot          & 221,9 MB                \\
        \bottomrule
    \end{tabular}
    \caption{Βασικά μεγέθη του CIC-IDS-2017 πριν και μετά την
    προεπεξεργασία.}\label{tab:dataset_sizes}
\end{table}

\section{Διάρθρωση της Αναφοράς}

Η υπόλοιπη αναφορά οργανώνεται ως εξής:

\begin{itemize}
    \item \textbf{Κεφάλαιο 2 - Θεωρητικό Υπόβαθρο.} Περιγράφονται οι
    θεμελιώδεις έννοιες που χρησιμοποιούνται σε όλα τα ερωτήματα:
    preprocessing, class imbalance, Pearson correlation,
    multicollinearity, κ.α.

    \item \textbf{Κεφάλαιο 3 - Ερώτημα 1: Exploratory Data Analysis.}
    Λεπτομερής περιγραφή της μεθοδολογίας και των αποτελεσμάτων του Q1.

    \item \textbf{Κεφάλαιο 4 - Ερώτημα 2: Ταξινόμηση.} Εκπαίδευση και
    αξιολόγηση μοντέλων (θα συμπληρωθεί).

    \item \textbf{Κεφάλαιο 5 - Ερώτημα 3: Συσταδοποίηση.} Unsupervised
    learning και εξωτερική αξιολόγηση (θα συμπληρωθεί).

    \item \textbf{Κεφάλαιο 6 - Συμπεράσματα και Μελλοντικές Κατευθύνσεις.}

    \item \textbf{Παράρτημα} με οδηγίες εγκατάστασης, εκτέλεσης και
    αναπαραγωγής.
\end{itemize}

%=============================================================================
% ΚΕΦΑΛΑΙΟ 2: ΘΕΩΡΗΤΙΚΟ ΥΠΟΒΑΘΡΟ
%=============================================================================
\chapter{Θεωρητικό Υπόβαθρο}

Σε αυτό το κεφάλαιο συνοψίζονται τα βασικά θεωρητικά εργαλεία που
χρησιμοποιούνται στα τρία ερωτήματα. Οι έννοιες παρουσιάζονται όπως εμφανίζονται
στο pipeline: από τον καθαρισμό δεδομένων (Q1) μέχρι την ταξινόμηση (Q2) και
τη συσταδοποίηση (Q3).

\section{Προεπεξεργασία Δεδομένων}

\subsection{Missing Values και Infinities}

Ένα πραγματικό σύνολο δεδομένων σπάνια είναι πλήρες. Στις δικτυακές
μετρήσεις του CIC-IDS-2017 οι \textit{rate} στήλες
(\texttt{Flow Bytes/s}, \texttt{Flow Packets/s}) υπολογίζονται ως διαιρέσεις
με τη \texttt{Flow Duration}. Όταν η διάρκεια μιας ροής είναι μηδέν
μικροδευτερόλεπτα (π.χ. single-packet flow), προκύπτουν τιμές
\(\pm\infty\). Η ανάγνωση αυτών από οποιονδήποτε
\texttt{StandardScaler}, \texttt{RobustScaler} ή από τον ίδιο το classifier
αποτυγχάνει.

Η τυπική πρακτική είναι:
\begin{enumerate}
    \item Αντικατάσταση \(\pm\infty\) με \texttt{NaN} (\texttt{df.replace}).
    \item Επιλογή μεταξύ \textbf{drop rows} ή \textbf{imputation}:
    \begin{itemize}
        \item \textbf{Drop}: απλό, διατηρεί την κατανομή, χάνει δεδομένα.
        \item \textbf{Median imputation}: ανθεκτικό σε outliers, αλλά
        εισάγει bias σε rate στήλες.
        \item \textbf{Zero imputation}: δηλώνει «μηδενικό throughput»,
        συνήθως λάθος σε flow-based μετρικές.
    \end{itemize}
\end{enumerate}

Όταν τα missing values είναι \(<1\%\) και σχετίζονται με ανώμαλες
παρατηρήσεις (π.χ. μηδενική διάρκεια), η drop-rows είναι συνήθως η
ασφαλέστερη επιλογή.

\subsection{Εξαιρετικά Διπλότυπα}

Τα exact duplicates δημιουργούν δύο προβλήματα:
\begin{itemize}
    \item \textbf{Πληθωρισμό κλάσεων}: αν χιλιάδες BENIGN γραμμές είναι
    πανομοιότυπες, το accuracy υπερεκτιμά τη γενικευτική ικανότητα του
    μοντέλου.
    \item \textbf{Data leakage μεταξύ train/test}: αν μια γραμμή εμφανιστεί
    και στα δύο split, το μοντέλο «απομνημονεύει» αντί να «μαθαίνει».
\end{itemize}

Γι' αυτό η αφαίρεση duplicates γίνεται \textbf{πριν} από κάθε split και
όχι αφού.

\section{Κλάσεις σε Ανισορροπία}

Όταν η αναλογία της μεγαλύτερης προς τη μικρότερη κλάση ξεπερνά το
\(\sim 10:1\), μιλάμε για \textbf{imbalanced classification}. Στο
CIC-IDS-2017 αυτή η αναλογία είναι \(\sim 195\,000:1\) (BENIGN vs.
Heartbleed), με τρεις άμεσες συνέπειες:

\begin{enumerate}
    \item \textbf{Παραπλανητική accuracy.} Ένας taxonomist «always
    BENIGN» επιτυγχάνει \(\sim 83\%\) χωρίς να ανιχνεύσει ούτε μία επίθεση.
    Γι' αυτό αναφέρουμε \textbf{macro-F1} και \textbf{per-class precision/recall}.

    \item \textbf{Προκατάληψη εκπαίδευσης.} Οι default αλγόριθμοι
    ελαχιστοποιούν το total loss, άρα ευνοούν τη majority class. Αντίμετρα:
    \begin{itemize}
        \item \texttt{class\_weight=\textquotesingle balanced\textquotesingle}: αντίστροφα βάρη
        στο loss.
        \item \textbf{SMOTE} (Synthetic Minority Over-sampling Technique):
        synthetic samples για τις minority classes.
        \item \textbf{Undersampling} της majority class.
    \end{itemize}

    \item \textbf{Stratified splitting υποχρεωτικό.} Ένα τυχαίο split
    μπορεί να στείλει όλες τις 11 Heartbleed γραμμές στο training set,
    αφήνοντας το test σύνολο χωρίς δείγματα για αξιολόγηση. Το
    \texttt{StratifiedShuffleSplit} διατηρεί τα αναλογικά ποσοστά ανά κλάση.
\end{enumerate}

\section{Pearson Correlation}

Ο \textbf{συντελεστής συσχέτισης Pearson} μεταξύ δύο μεταβλητών \(X, Y\)
ορίζεται ως:
\begin{equation}
    r_{XY} = \frac{\mathrm{Cov}(X, Y)}{\sigma_X \sigma_Y}
           = \frac{\sum_{i=1}^{n}(x_i-\bar x)(y_i-\bar y)}
                  {\sqrt{\sum_{i=1}^{n}(x_i-\bar x)^2}\,
                   \sqrt{\sum_{i=1}^{n}(y_i-\bar y)^2}}
\end{equation}
με \(r \in [-1, +1]\). Τιμές \(\pm 1\) δηλώνουν τέλεια γραμμική εξάρτηση,
\(r \approx 0\) καμία \emph{γραμμική} εξάρτηση (αλλά μπορεί να υπάρχει
μη-γραμμική).

\subsection{Τυφλά Σημεία}

Η Pearson έχει τρεις γνωστούς περιορισμούς που είναι χρήσιμοι στην ερμηνεία
των αποτελεσμάτων:
\begin{itemize}
    \item Δεν ανιχνεύει \textbf{μη-γραμμικές} σχέσεις (π.χ.
    \(Y = X^2\) δίνει \(r \approx 0\)).
    \item Είναι ιδιαίτερα ευαίσθητη σε \textbf{outliers}, ειδικά σε
    κατανομές heavy-tailed.
    \item Απαιτεί \textbf{αριθμητικές} μεταβλητές. Το label encoding που
    εφαρμόζουμε στο \texttt{Label} επιβάλλει αυθαίρετη διάταξη, οπότε
    η \texttt{corrwith(Label)} είναι μόνο \textit{ενδεικτική}. Για
    αυστηρή αξιολόγηση feature importance χρησιμοποιείται
    \textbf{mutual information} ή tree-based importance.
\end{itemize}

\section{Multicollinearity}

Δύο features με \(|r| \to 1\) μεταφέρουν την ίδια πληροφορία. Η
\textbf{πολυσυγγραμμικότητα} (multicollinearity) βλάπτει τα μοντέλα με τρεις
τρόπους:

\begin{enumerate}
    \item \textbf{Γραμμικά μοντέλα} (Logistic Regression): οι συντελεστές
    γίνονται ασταθείς. Με δύο σχεδόν πανομοιότυπα features, ο αλγόριθμος
    μπορεί να δώσει \(w_1=+10,\ w_2=-10\) ή \(w_1=+5,\ w_2=-5\),
    μαθηματικά ισοδύναμα αλλά ανερμήνευτα. Η τυπική απόκλιση των
    coefficients εκρήγνυται.
    \item \textbf{Αλγόριθμοι βασισμένοι σε αποστάσεις} (K-Means, KNN,
    DBSCAN): η επανειλημμένη διάσταση μετράει διπλά στην ευκλείδεια
    απόσταση, οπότε αυτή η διάσταση κυριαρχεί τεχνητά.
    \item \textbf{Υπολογιστικά}: περισσότερες στήλες σημαίνει πιο αργή
    εκπαίδευση χωρίς κέρδος πληροφορίας.
\end{enumerate}

Η στρατηγική που ακολουθούμε είναι να κρατήσουμε μία μόνο στήλη ανά
cluster υψηλής συσχέτισης (με threshold \(|r|>0.95\)), διατηρώντας την
πρώτη στη σειρά των στηλών και πετώντας τις άλλες.

\section{log1p Transform και Skewness}

Πολλά features του CIC-IDS-2017 εκτείνονται σε πολλές τάξεις μεγέθους
(π.χ. \texttt{Flow Duration} από μs έως λεπτά) και έχουν skewness
\(>100\) με kurtosis στις χιλιάδες. Χωρίς μετασχηματισμό, τα ιστογράμματα
συμπτύσσονται σε μία κορυφή κοντά στο 0 με λεπτή ουρά και οι boxplots
είναι οπτικά άχρηστα.

Ο μετασχηματισμός \textbf{log1p}:
\begin{equation}
    \texttt{log1p}(x) = \log(1+x)
\end{equation}
έχει τρία πλεονεκτήματα:
\begin{itemize}
    \item Είναι ασφαλής για \(x=0\) (ενώ \(\log(0) \to -\infty\)).
    \item Για μεγάλα \(x\), \(\texttt{log1p}(x) \approx \log(x)\).
    \item Για μικρά \(x\), η Taylor expansion δίνει \(\texttt{log1p}(x) \approx x\).
\end{itemize}

\section{Stratified Sampling για Οπτικοποίηση}

Για ένα dataset με \(\sim 2.57\) εκατομμύρια γραμμές, το KDE και τα
heatmaps γίνονται απαγορευτικά αργά. Η στρατηγική μας είναι:
\begin{itemize}
    \item \textbf{Στατιστικά} (\texttt{describe}, correlation matrix):
    ολόκληρο το dataset. Δεν υπάρχει λόγος να θυσιάσουμε ακρίβεια.
    \item \textbf{Οπτικοποιήσεις}: stratified sample \(50\,000\) γραμμών με
    \texttt{StratifiedShuffleSplit}, ώστε κάθε κλάση (ακόμα και η
    Heartbleed με 11 γραμμές) να εκπροσωπείται.
\end{itemize}

\section{Αξιολόγηση Ταξινομητών (Q2)}

\subsection{Confusion Matrix}

Για multi-class classification με \(C\) κλάσεις, η confusion matrix
\(M \in \mathbb{N}^{C \times C}\) ορίζεται ως
\(M_{ij} = \) πλήθος δειγμάτων με πραγματική κλάση \(i\) που ταξινομήθηκαν
στην κλάση \(j\). Η διαγώνιος είναι τα σωστά (TP ανά κλάση).

\subsection{Precision, Recall, F1}

Για κάθε κλάση \(c\):
\begin{align}
    \text{Precision}_c &= \frac{TP_c}{TP_c + FP_c}, \\
    \text{Recall}_c    &= \frac{TP_c}{TP_c + FN_c}, \\
    F1_c               &= 2 \cdot \frac{\text{Precision}_c \cdot \text{Recall}_c}
                                        {\text{Precision}_c + \text{Recall}_c}.
\end{align}

Το \textbf{macro-F1} είναι ο αστάθμιστος μέσος όρος των \(F1_c\) -
κάθε κλάση μετράει το ίδιο, ανεξάρτητα από το μέγεθός της. Το
\textbf{weighted-F1} σταθμίζει με το support, επομένως υπερεκτιμά τα
μοντέλα που τα πάνε καλά στη majority class.

\subsection{Grid Search Cross-Validation}

Το \texttt{GridSearchCV} εξαντλητικά δοκιμάζει όλους τους συνδυασμούς
hyperparameters σε \(k\)-fold cross-validation (συνήθως \(k=5\)). Η
\textbf{stratified} παραλλαγή διατηρεί την κατανομή των κλάσεων σε κάθε
fold, κρίσιμο για imbalanced data.

\section{Συσταδοποίηση (Q3)}

\subsection{K-Means}

Ο αλγόριθμος K-Means ελαχιστοποιεί την εντός-συστάδας διασπορά (inertia):
\begin{equation}
    \min_{\{C_k\}, \{\mu_k\}} \sum_{k=1}^{K} \sum_{x \in C_k} \|x - \mu_k\|^2.
\end{equation}

Επιλογή του \(K\): elbow method (inertia vs \(K\)) ή silhouette score.

\subsection{DBSCAN}

Σε αντίθεση με το K-Means, το DBSCAN δεν απαιτεί προκαθορισμένο \(K\),
βρίσκει συστάδες αυθαίρετου σχήματος και χαρακτηρίζει σημεία ως
\textbf{noise} αν δεν ανήκουν σε κάποια πυκνή περιοχή. Ορίζεται από δύο
παραμέτρους:
\begin{itemize}
    \item \(\varepsilon\) (eps): ακτίνα γειτονιάς.
    \item \texttt{min\_samples}: ελάχιστα σημεία σε μια πυκνή περιοχή.
\end{itemize}

\subsection{Hierarchical Clustering}

Κατασκευάζει ένα δενδρόγραμμα (dendrogram) συγχωνεύοντας επαναληπτικά τα
πιο κοντινά σημεία/συστάδες. Η linkage function (single, complete,
average, Ward) καθορίζει πώς μετράται η απόσταση μεταξύ συστάδων.

\subsection{External Validation με ARI}

Όταν οι \textit{πραγματικές} ετικέτες είναι γνωστές (όπως εδώ), μπορούμε
να αξιολογήσουμε τις συστάδες με \textbf{Adjusted Rand Index (ARI)} ή
\textbf{Normalized Mutual Information (NMI)}. Αυτό μας λέει πόσο καλά η
δομή των δεδομένων αντιστοιχεί στις κλάσεις επίθεσης, χωρίς να έχει
χρησιμοποιηθεί καμία ετικέτα κατά την εκπαίδευση.

%=============================================================================
% ΚΕΦΑΛΑΙΟ 3: ΕΡΩΤΗΜΑ 1 - EDA
%=============================================================================
\chapter{Ερώτημα 1: Exploratory Data Analysis}

\section{Επισκόπηση του Pipeline}

Το Q1 υλοποιήθηκε ως ένα μοναδικό Jupyter notebook
(\texttt{notebooks/q1\_eda.ipynb}), εύκολα επανεκτελέσιμο μέσω
\texttt{nbconvert}. Το pipeline που ακολουθεί είναι αυτό:

\begin{enumerate}
    \item Φόρτωση όλων των 8 CSV αρχείων και συγχώνευσή τους.
    \item Κανονικοποίηση ονομάτων στηλών και ετικετών.
    \item Αντιμετώπιση Inf και NaN.
    \item Αφαίρεση duplicates.
    \item Περιγραφικά στατιστικά και οπτικοποίηση κατανομής κλάσεων.
    \item Ιστογράμματα, boxplots, correlation heatmap.
    \item Feature selection (zero-variance + high-correlation).
    \item Αποθήκευση σε Parquet για επαναχρησιμοποίηση από Q2/Q3.
\end{enumerate}

\section{Φόρτωση και Επισκόπηση}

\subsection{Ενοποίηση των 8 CSV}

Η βοηθητική συνάρτηση \texttt{load\_cic\_ids\_2017} (αρχείο
\texttt{utils/helpers.py}) διαβάζει ένα προς ένα τα 8 CSV αρχεία με
\texttt{encoding=\textquotesingle latin1\textquotesingle} για να
αντιμετωπίσει τα corrupted bytes στα labels των Web Attacks, κάνει
concatenation και κρατά προέλευση μέσω της στήλης \texttt{\_\_source\_file}
(αφαιρείται αργότερα).

\begin{tcolorbox}[myCustomStyle={Φόρτωση Dataset}]
\begin{lstlisting}[style=myPythonStyle]
from utils.helpers import load_cic_ids_2017, RANDOM_STATE
import numpy as np, pandas as pd

df = load_cic_ids_2017()
print(df.shape)   # (2830743, 80)
\end{lstlisting}
\end{tcolorbox}

Μέσα στη \texttt{\_standardize\_columns} γίνεται και ένα κρίσιμο βήμα: η
στήλη \texttt{Fwd Header Length.1} - ένα αντίγραφο της
\texttt{Fwd Header Length} που εξάγει λανθασμένα το CICFlowMeter -
αφαιρείται αυτόματα.

\subsection{Καθαρισμός Ετικετών}

Οι ετικέτες των Web Attacks στα raw CSV περιέχουν το χαρακτήρα U+FFFD (το
Unicode replacement character) στη θέση ενός en-dash:
\texttt{"Web Attack \(\square\) Brute Force"}. Η \texttt{\_clean\_labels}
τα κανονικοποιεί σε \texttt{"Web Attack Brute Force"}, \texttt{"Web Attack XSS"},
\texttt{"Web Attack Sql Injection"}.

\section{Καθαρισμός Δεδομένων}

\subsection{Infinities και Missing Values}

Οι στήλες \texttt{Flow Bytes/s} και \texttt{Flow Packets/s} περιέχουν
\(4\,376\) τιμές \(\pm\infty\) (από διαίρεση με \texttt{Flow Duration}=0).
Ακολουθούμε τη στρατηγική replace-\&-drop:

\begin{tcolorbox}[myCustomStyle={Inf και NaN}]
\begin{lstlisting}[style=myPythonStyle]
from utils.helpers import handle_infinities

df = handle_infinities(df)   # (+/-inf) -> NaN
print(df.isna().sum().sum()) # 4376 NaN total
before = len(df)
df = df.dropna().reset_index(drop=True)
print(before - len(df), "rows dropped")   # 2867 rows
\end{lstlisting}
\end{tcolorbox}

Το αποτέλεσμα είναι \textbf{2.867} γραμμές που αφαιρέθηκαν, δηλαδή
\(0.101\%\) του dataset - ποσοστό πολύ μικρό για να επηρεάσει τη δομή.

\subsection{Duplicates}

Η \texttt{drop\_duplicates\_report} αναφέρει \textbf{255.236} exact
duplicates (\(\sim 9\%\) μετά το NaN handling). Αυτά αφαιρούνται \textit{πριν}
το οποιοδήποτε split, ώστε το test set του Q2 να είναι πραγματικά unseen.

\begin{table}[H]
    \centering
    \begin{tabular}{lrr}
        \toprule
        \textbf{Στάδιο} & \textbf{Γραμμές} & \textbf{Διαφορά}       \\
        \midrule
        Αρχικό           & 2.830.743        & ---                    \\
        Μετά drop NaN    & 2.827.876        & \(-2\,867\)            \\
        Μετά drop dup    & 2.572.640        & \(-255\,236\)          \\
        \bottomrule
    \end{tabular}
    \caption{Προοδευτική μείωση του μεγέθους του dataset κατά τον καθαρισμό.}
    \label{tab:cleaning_stages}
\end{table}

\section{Κατανομή Κλάσεων}

Η κατανομή των 15 κλάσεων φαίνεται στον Πίνακα~\ref{tab:label_distribution}
και οπτικοποιείται στο Σχήμα~\ref{fig:label_distribution}.

\begin{table}[H]
    \centering
    \small
    \begin{tabular}{lrr}
        \toprule
        \textbf{Κλάση}                & \textbf{Πλήθος} & \textbf{Ποσοστό} \\
        \midrule
        BENIGN                        & 2.146.899        & 83,45\%          \\
        DoS Hulk                      & 172.846          & 6,72\%           \\
        DDoS                          & 128.014          & 4,98\%           \\
        PortScan                      & 90.694           & 3,53\%           \\
        DoS GoldenEye                 & 10.286           & 0,40\%           \\
        FTP-Patator                   & 5.931            & 0,23\%           \\
        DoS slowloris                 & 5.385            & 0,21\%           \\
        DoS Slowhttptest              & 5.228            & 0,20\%           \\
        SSH-Patator                   & 3.219            & 0,13\%           \\
        Bot                           & 1.948            & 0,08\%           \\
        Web Attack Brute Force        & 1.470            & 0,06\%           \\
        Web Attack XSS                & 652              & 0,03\%           \\
        Infiltration                  & 36               & 0,0014\%         \\
        Web Attack Sql Injection      & 21               & 0,0008\%         \\
        Heartbleed                    & 11               & 0,0004\%         \\
        \midrule
        \textbf{Σύνολο}               & \textbf{2.572.640} & \textbf{100\%}  \\
        \bottomrule
    \end{tabular}
    \caption{Κατανομή των 15 κλάσεων μετά τον καθαρισμό. Το imbalance
    ratio (max/min) είναι \(\approx 195\,173\times\).}
    \label{tab:label_distribution}
\end{table}

\begin{figure}[H]
    \centering
    \begin{subfigure}[b]{0.48\textwidth}
        \includegraphics[width=\textwidth]{q1_label_distribution_bar.png}
        \caption{Λογαριθμικό bar chart}
    \end{subfigure}
    \hfill
    \begin{subfigure}[b]{0.48\textwidth}
        \includegraphics[width=\textwidth]{q1_label_distribution_pie.png}
        \caption{Pie chart (επικρατεί η BENIGN)}
    \end{subfigure}
    \caption{Κατανομή κλάσεων στο καθαρισμένο CIC-IDS-2017. Το αριστερό
    σχήμα χρησιμοποιεί λογαριθμικό άξονα ώστε να φαίνονται και οι
    σπάνιες κλάσεις (Heartbleed, Sql Injection, Infiltration).}
    \label{fig:label_distribution}
\end{figure}

\textbf{Παρατηρήσεις.} Η ακραία ανισορροπία (η \texttt{BENIGN} είναι
\(>195\,000\times\) πιο συχνή από την \texttt{Heartbleed}) καθορίζει τις
επιλογές του Q2: θα απαιτηθεί \texttt{class\_weight="balanced"},
macro-F1 ως κύρια metric και stratified splitting.

\section{Περιγραφικά Στατιστικά}

Για κάθε αριθμητική στήλη υπολογίστηκαν mean, median, std, min, 25\%, 50\%,
75\%, max, skewness και kurtosis. Τα πλήρη αποτελέσματα αποθηκεύτηκαν στο
\texttt{outputs/results/q1\_describe.csv}. Τα βασικά ευρήματα:

\begin{itemize}
    \item Πολλές στήλες έχουν \textbf{skewness} \(>100\) και
    \textbf{kurtosis} στις χιλιάδες - επιβεβαιώνοντας κατανομές
    heavy-tailed που απαιτούν log-transform για οπτικοποίηση.
    \item Οι στήλες \texttt{Fwd Avg Bytes/Bulk},
    \texttt{Bwd Avg Packets/Bulk} κ.ά. έχουν \texttt{std}=0, δηλαδή είναι
    \textbf{σταθερές} σε όλες τις \(\sim 2.57\)M γραμμές και θα αφαιρεθούν
    στο βήμα feature selection.
\end{itemize}

\section{Οπτικοποιήσεις}

\subsection{Ιστογράμματα Top Features}

Επιλέχθηκαν 20 features με τη μεγαλύτερη διασπορά για plot. Για λόγους
ευκρίνειας ο άξονας \(x\) είναι σε log scale (\texttt{log1p}) και το KDE
υπολογίστηκε σε stratified sample \(50\,000\) γραμμών.

\begin{figure}[H]
    \centering
    \includegraphics[width=0.95\textwidth]{q1_histograms_kde.png}
    \caption{Ιστογράμματα με KDE για τις 20 στήλες μεγαλύτερης διασποράς
    (log1p scale). Οι περισσότερες κατανομές είναι έντονα λοξές και
    heavy-tailed.}\label{fig:histograms}
\end{figure}

\subsection{Boxplots}

Οι boxplots (Σχήμα~\ref{fig:boxplots}) αποκαλύπτουν έντονα «νέφοι»
outliers στις rate και duration στήλες. Αυτά δεν είναι «θόρυβος» προς
αφαίρεση - είναι ακριβώς η υπογραφή των DDoS/DoS επιθέσεων (πολύ υψηλά
throughput, πολύ μακριές διάρκειες).

\begin{figure}[H]
    \centering
    \includegraphics[width=0.95\textwidth]{q1_boxplots_top20.png}
    \caption{Boxplots (log1p scale) για τις 20 στήλες μεγαλύτερης
    διασποράς. Τα outliers πάνω από τα whiskers είναι διαγνωστικά για
    επιθέσεις DoS/DDoS και δεν αφαιρούνται.}\label{fig:boxplots}
\end{figure}

\subsection{Correlation Heatmap}

Η Pearson correlation matrix (77 \(\times\) 77) οπτικοποιείται ως
heatmap (Σχήμα~\ref{fig:corr_heatmap}). Τα βαθυκόκκινα και βαθυγάλανα
τετράγωνα είναι τα ζεύγη που θα στοχεύσουμε στο feature selection.

\begin{figure}[H]
    \centering
    \includegraphics[width=0.95\textwidth]{q1_correlation_heatmap.png}
    \caption{Pearson correlation heatmap των 77 αριθμητικών features.
    Τα μεγάλα κόκκινα clusters αντιστοιχούν σε αλγοριθμικές ταυτότητες
    του CICFlowMeter (π.χ. \texttt{Subflow Fwd Bytes} \(\approx\)
    \texttt{Total Length of Fwd Packets}).}\label{fig:corr_heatmap}
\end{figure}

\subsection{Top Correlated Pairs}

Τα 15 ζεύγη με τη μεγαλύτερη \(|r|\) φαίνονται στον
Πίνακα~\ref{tab:top_correlated}. Τα επτά πρώτα ζεύγη έχουν τέλεια
συσχέτιση (\(r=1.0\)) - δηλαδή είναι \textbf{αλγοριθμικές ταυτότητες}
που ο CICFlowMeter εξάγει δύο φορές με διαφορετικά ονόματα.

\begin{table}[H]
    \centering
    \small
    \begin{tabular}{llr}
        \toprule
        \textbf{Feature A}            & \textbf{Feature B}           & \textbf{\(r\)} \\
        \midrule
        Bwd Packet Length Mean        & Avg Bwd Segment Size         & 1,0000         \\
        Total Fwd Packets             & Subflow Fwd Packets          & 1,0000         \\
        Total Backward Packets        & Subflow Bwd Packets          & 1,0000         \\
        Fwd URG Flags                 & CWE Flag Count               & 1,0000         \\
        Fwd Packet Length Mean        & Avg Fwd Segment Size         & 1,0000         \\
        Fwd PSH Flags                 & SYN Flag Count               & 1,0000         \\
        Total Length of Bwd Packets   & Subflow Bwd Bytes            & 1,0000         \\
        Total Length of Fwd Packets   & Subflow Fwd Bytes            & 1,0000         \\
        Total Backward Packets        & Subflow Fwd Packets          & 0,9991         \\
        Subflow Fwd Packets           & Subflow Bwd Packets          & 0,9991         \\
        Total Fwd Packets             & Subflow Bwd Packets          & 0,9991         \\
        Total Fwd Packets             & Total Backward Packets       & 0,9991         \\
        Flow Duration                 & Fwd IAT Total                & 0,9985         \\
        Flow IAT Max                  & Fwd IAT Max                  & 0,9981         \\
        Packet Length Mean            & Average Packet Size          & 0,9978         \\
        \bottomrule
    \end{tabular}
    \caption{Top-15 ζεύγη με \(|r| > 0.997\). Τα 8 πρώτα είναι
    αλγοριθμικές ταυτότητες του CICFlowMeter.}\label{tab:top_correlated}
\end{table}

\subsection{Correlation με το Label}

Καμία στήλη δεν έχει \(|r| > 0.31\) με το label (εξαρτώμενο από label
encoding), επιβεβαιώνοντας ότι το πρόβλημα είναι γνήσια \emph{πολυμεταβλητό}.
Οι 5 πιο σχετικές στήλες είναι \texttt{Fwd IAT Std} (0,30),
\texttt{Idle Mean} (0,28), \texttt{Bwd Packet Length Std} (0,28),
\texttt{Idle Max} (0,28), \texttt{Idle Min} (0,28).

\textbf{Σημείωση:} Τέσσερα από τα παραπάνω πέντε features (\texttt{Idle
Mean/Max/Min}, \texttt{Bwd Packet Length Std}) θα αφαιρεθούν στην επόμενη
ενότητα ως υψηλά συσχετισμένα. Αυτό \emph{δεν} αφαιρεί πληροφορία: ανήκουν
στο ίδιο cluster με features που \emph{διατηρούνται} (π.χ.
\texttt{Flow IAT Max} με \(r=0{,}276\) προς το label, που εκπροσωπεί τη
ομάδα Idle/IAT-Max). Έτσι ο επόμενος ταξινομητής βλέπει την ίδια διακριτική
σήμανση μέσω του εκπροσώπου της κάθε ομάδας.

\section{Feature Selection}

\subsection{Κριτήρια Αφαίρεσης}

Εφαρμόσαμε τρία κριτήρια, με την ακόλουθη σειρά:

\begin{table}[H]
    \centering
    \small
    \begin{tabular}{p{4cm}p{8cm}}
        \toprule
        \textbf{Κριτήριο}                    & \textbf{Αιτιολογία}                                \\
        \midrule
        Bookkeeping                          & Βοηθητικές στήλες (\texttt{\_\_source\_file})      \\
                                             & που δεν είναι feature.                             \\
        \midrule
        Zero variance (\(\sigma^2=0\))       & Σταθερές στήλες: καμία πληροφορία,                 \\
                                             & καταρρίπτουν κάθε scaler/model.                    \\
        \midrule
        High correlation (\(|r|>0.95\))      & Multicollinear near-duplicates:                    \\
                                             & κρατάμε μία ανά cluster.                           \\
        \bottomrule
    \end{tabular}
    \caption{Τα τρία κριτήρια feature selection και η αιτιολογία τους.}
    \label{tab:selection_criteria}
\end{table}

\subsection{Αποτελέσματα}

Συνολικά αφαιρέθηκαν \textbf{31 στήλες} (Πίνακας~\ref{tab:removed_features}):

\begin{itemize}
    \item \textbf{1} bookkeeping (\texttt{\_\_source\_file}).
    \item \textbf{8} zero-variance: όλες οι Bulk-rate (\texttt{Fwd/Bwd Avg
    Bytes/Bulk}, \texttt{.../Packets/Bulk}, \texttt{.../Bulk Rate}) και
    \texttt{Bwd PSH Flags}, \texttt{Bwd URG Flags}. Όλες είναι μηδενικές
    σε όλα τα \(\sim 2.57\)M flows.
    \item \textbf{22} highly correlated: αλγοριθμικές ταυτότητες του
    CICFlowMeter και near-duplicates (\texttt{Average Packet Size},
    \texttt{Avg Bwd Segment Size}, \texttt{Subflow Fwd Packets},
    κλπ.).
\end{itemize}

\begin{table}[H]
    \centering
    \scriptsize
    \begin{tabular}{p{4.5cm}p{7.5cm}}
        \toprule
        \textbf{Στήλη}                  & \textbf{Λόγος}                                                  \\
        \midrule
        \_\_source\_file                & bookkeeping / non-feature helper column                         \\
        Bwd PSH Flags                   & zero variance (constant)                                        \\
        Bwd URG Flags                   & zero variance (constant)                                        \\
        Fwd Avg Bytes/Bulk              & zero variance (constant)                                        \\
        Fwd Avg Packets/Bulk            & zero variance (constant)                                        \\
        Fwd Avg Bulk Rate               & zero variance (constant)                                        \\
        Bwd Avg Bytes/Bulk              & zero variance (constant)                                        \\
        Bwd Avg Packets/Bulk            & zero variance (constant)                                        \\
        Bwd Avg Bulk Rate               & zero variance (constant)                                        \\
        Average Packet Size             & highly correlated (\(|r|>0.95\)) with a kept feature            \\
        Avg Bwd Segment Size            & highly correlated (\(|r|>0.95\))                                \\
        Avg Fwd Segment Size            & highly correlated (\(|r|>0.95\))                                \\
        Bwd Packet Length Mean          & highly correlated (\(|r|>0.95\))                                \\
        Bwd Packet Length Std           & highly correlated (\(|r|>0.95\))                                \\
        CWE Flag Count                  & highly correlated (\(|r|>0.95\))                                \\
        ECE Flag Count                  & highly correlated (\(|r|>0.95\))                                \\
        Fwd IAT Max                     & highly correlated (\(|r|>0.95\))                                \\
        Fwd IAT Total                   & highly correlated (\(|r|>0.95\))                                \\
        Fwd Packet Length Std           & highly correlated (\(|r|>0.95\))                                \\
        Fwd Packets/s                   & highly correlated (\(|r|>0.95\))                                \\
        Idle Max                        & highly correlated (\(|r|>0.95\))                                \\
        Idle Mean                       & highly correlated (\(|r|>0.95\))                                \\
        Idle Min                        & highly correlated (\(|r|>0.95\))                                \\
        Packet Length Std               & highly correlated (\(|r|>0.95\))                                \\
        SYN Flag Count                  & highly correlated (\(|r|>0.95\))                                \\
        Subflow Bwd Bytes               & highly correlated (\(|r|>0.95\))                                \\
        Subflow Bwd Packets             & highly correlated (\(|r|>0.95\))                                \\
        Subflow Fwd Bytes               & highly correlated (\(|r|>0.95\))                                \\
        Subflow Fwd Packets             & highly correlated (\(|r|>0.95\))                                \\
        Total Backward Packets          & highly correlated (\(|r|>0.95\))                                \\
        Total Length of Bwd Packets     & highly correlated (\(|r|>0.95\))                                \\
        \bottomrule
    \end{tabular}
    \caption{Οι 31 στήλες που αφαιρέθηκαν στο feature
    selection.}\label{tab:removed_features}
\end{table}

\subsection{Τελική Λίστα Features}

Μετά το feature selection έχουμε \textbf{47 features + Label} (48 στήλες
συνολικά), οι οποίες αποθηκεύτηκαν στο
\texttt{outputs/results/q1\_final\_feature\_list.csv}. Ενδεικτικά, οι πρώτες
δέκα: \texttt{Destination Port}, \texttt{Flow Duration},
\texttt{Total Fwd Packets}, \texttt{Total Length of Fwd Packets},
\texttt{Fwd Packet Length Max}, \texttt{Fwd Packet Length Min},
\texttt{Fwd Packet Length Mean}, \texttt{Bwd Packet Length Max},
\texttt{Bwd Packet Length Min}, \texttt{Flow Bytes/s}.

\section{Αποθήκευση του Clean Snapshot}

Το καθαρισμένο dataset σώζεται σε Parquet
(\texttt{data/cic\_ids\_2017\_clean.parquet}, 221,9 MB) μέσω της
\texttt{save\_clean\_cached}. Τα Q2 και Q3 το φορτώνουν άμεσα με
\texttt{load\_clean\_cached()} χωρίς να επαναλάβουν το pipeline.

\textbf{Γιατί Parquet και όχι CSV:}
\begin{itemize}
    \item \textbf{Columnar storage}: αν το Q3 θέλει μόνο 10 στήλες, τις
    διαβάζει χωρίς να φορτώσει τις υπόλοιπες 38.
    \item \textbf{Compressed}: 221,9 MB αντί για \(\sim\)900 MB σε CSV.
    \item \textbf{Type-preserved}: διατηρεί τα int32/float32 που
    προσεκτικά ορίσαμε στη \texttt{\_optimize\_memory}.
    \item \textbf{Fast I/O}: τα Q2/Q3 φορτώνουν σε \(<2\) sec αντί για
    \(\sim 60\) sec CSV parsing.
\end{itemize}

\section{Σύνοψη Q1}

\begin{table}[H]
    \centering
    \begin{tabular}{lr}
        \toprule
        \textbf{Μέγεθος}                            & \textbf{Τιμή}        \\
        \midrule
        Αρχικές γραμμές                             & 2.830.743            \\
        Rows με \(\pm\infty\) \(\to\) NaN           & 4.376 (τιμές)        \\
        Rows που έπεσαν λόγω NaN                    & 2.867 (\(0,101\%\))  \\
        Exact duplicates που αφαιρέθηκαν            & 255.236 (\(\sim 9\%\)) \\
        Τελικές γραμμές                             & 2.572.640            \\
        Αρχικές στήλες                              & 79 + Label           \\
        Στήλες που αφαιρέθηκαν                      & 31                   \\
        Τελικές στήλες                              & 47 + Label           \\
        Κλάσεις                                     & 15                   \\
        Imbalance ratio (max/min)                   & \(\sim 195\,173\times\) \\
        Parquet snapshot                            & 221,9 MB             \\
        \bottomrule
    \end{tabular}
    \caption{Βασικά αποτελέσματα του Q1.}\label{tab:q1_summary}
\end{table}

%=============================================================================
% ΚΕΦΑΛΑΙΟ 4: ΕΡΩΤΗΜΑ 2 - CLASSIFICATION
%=============================================================================
\chapter{Ερώτημα 2: Ταξινόμηση}\label{ch:q2}

Στο κεφάλαιο αυτό εκπαιδεύουμε τρεις ταξινομητές (Logistic Regression,
Decision Tree, Random Forest) στο καθαρισμένο dataset του Q1 και τους
αξιολογούμε με δύο διαφορετικά πρωτόκολλα hyperparameter selection:
(α) \emph{Σενάριο A} - held-out validation set σε stratified split
60/20/20, και (β) \emph{Σενάριο B} - 5-fold Stratified CV μέσω
\texttt{GridSearchCV} πάνω στο 80\% training, με το ίδιο test set σε
όλες τις παραλλαγές.

\section{Μεθοδολογία}\label{sec:q2_methodology}

\subsection{Stratified Subsampling με Per-Class Cap}

Το πλήρες cleaned dataset έχει 2.572.640 γραμμές. Για να γίνει
υπολογιστικά εφικτή η εξαντλητική grid search (ιδίως για το Random
Forest με 18 combinations \(\times\) 5 CV folds), εφαρμόζουμε
\textbf{stratified subsampling με per-class cap} στις
\textsc{25.000} rows. Οι μειονοτικές κλάσεις (Heartbleed \(=11\),
SQL~Injection \(=21\), Infiltration \(=36\), \dots) διατηρούνται
\textbf{πλήρως} - το cap ενεργοποιείται μόνο για τις τέσσερις τεράστιες
(\texttt{BENIGN}, \texttt{DoS~Hulk}, \texttt{DDoS}, \texttt{PortScan}).

\textbf{Συνέπεια στο imbalance.} Το imbalance ratio πέφτει από
\(195\,173{:}1\) στο raw data σε \(\sim 2\,272{:}1\)
(25.000 BENIGN έναντι 11 Heartbleed). Παραμένει αρκετά σοβαρό ώστε το
\texttt{class\_weight="balanced"} να διατηρεί το νόημά του· η
παρακολούθηση macro-F1 συνεχίζει να είναι απαραίτητη για να
αντιληφθούμε την ποιότητα στις rare classes.

\textbf{Τελικό μέγεθος}: 134.187 rows \(\times\) 47 features + Label.

\textbf{Configurable mode.} Στο notebook υπάρχει ορισμένη μεταβλητή
\texttt{SUBSAMPLE\_CAP} με τρία πρακτικά presets: \texttt{25\_000}
(fast, \(\sim 134\)K rows, \(\sim 15\text{-}25\) min - default που
αναφέρεται εδώ), \texttt{100\_000} (medium, \(\sim 500\)K rows,
\(\sim 2\text{-}4\) hours) και \texttt{None} (full dataset,
\(\sim 2{,}5\)M rows, \(\sim 12+\) hours σε σύγχρονο laptop). Ο
κύριος bottleneck στο full mode είναι το LR με \texttt{saga} solver,
που μπορεί να χρειαστεί 1-3 ώρες ανά grid combination.

\subsection{Encoding, Scaling και Imbalance Handling}

\begin{itemize}
    \item \textbf{LabelEncoder} στο target: μετατρέπει τα 15 string
    labels σε ακεραίους \([0, 14]\), όπως περιμένουν οι multi-class
    classifiers του scikit-learn. One-hot δεν έχει νόημα στο target
    (θα γινόταν multi-output regression).
    \item \textbf{StandardScaler μέσα σε \texttt{Pipeline}} για το
    Logistic Regression. Το pipeline εξασφαλίζει ότι ο scaler γίνεται
    \texttt{fit} \emph{μόνο} στο training fold σε κάθε CV split -
    αποφεύγοντας data leakage. Για Decision Tree και Random Forest
    παραλείπεται: τα splits του δέντρου είναι invariant σε monotonic
    transforms του feature.
    \item \textbf{\texttt{class\_weight="balanced"}} σε κάθε αλγόριθμο.
    Ζυγίζει κάθε row με \(n/(C\cdot n_c)\) όπου \(n\) το σύνολο, \(C\)
    ο αριθμός κλάσεων και \(n_c\) η συχνότητα της κλάσης. Χωρίς αυτό,
    ένας classifier θα έβγαζε \(\sim 80\%\) accuracy απλώς
    προβλέποντας πάντα \texttt{BENIGN}. Εναλλακτικές (SMOTE,
    undersampling) εγκυμονούν data-leakage risks αν δεν εφαρμοστούν
    per-fold.
\end{itemize}

\subsection{Πρωτόκολλο Splits}

Το subsampled dataset σπάει σε δύο βήματα:
\begin{enumerate}
    \item Stratified 80/20 \(\to\) \texttt{trainval} / \texttt{test}.
    Το \texttt{test} (26.838 rows) είναι ο \textbf{τυφλός μάρτυρας}:
    αγγίζεται \emph{μόνο} για το final reporting και των δύο σεναρίων.
    \item Από το \texttt{trainval} σπάμε ξανά stratified 75/25 \(\to\)
    \texttt{train} (80.511 rows) / \texttt{val} (26.838 rows). Αυτό
    δίνει τον τελικό 60/20/20 λόγο που ζητά η εκφώνηση.
\end{enumerate}

\begin{itemize}
    \item \textbf{Σενάριο A} (split-based grid search): για κάθε
    combination στο grid, fit σε \texttt{train}, score σε
    \texttt{val} με F1-weighted. Επιλέγουμε best params, refit στο
    \texttt{trainval}, evaluate στο \texttt{test}.
    \item \textbf{Σενάριο B} (5-fold StratifiedKFold μέσω
    \texttt{GridSearchCV}) πάνω στο ίδιο \texttt{trainval} και ίδιο
    grid. Best params \(\to\) refit στο \texttt{trainval}, evaluate
    στο \texttt{test}.
\end{itemize}

Το κοινό \texttt{test} σε κάθε παραλλαγή είναι αυτό που καθιστά
δυνατή τη μεταξύ τους σύγκριση.

\subsection{Scoring Metric για την Επιλογή Hyperparameters}

Χρησιμοποιούμε \textbf{F1-weighted} και στα δύο σενάρια. Στο final
reporting εμφανίζουμε \emph{και} F1-macro. Λόγοι:
\begin{itemize}
    \item \textbf{Accuracy} εξαπατά σε imbalanced multi-class (\(\sim
    80\%\) απλώς από το να προβλέπει BENIGN πάντα).
    \item \textbf{F1-macro} δίνει ίσο βάρος σε κάθε κλάση, αλλά οι
    rare κλάσεις (Heartbleed \(=11\) rows) έχουν τέτοια variance ώστε
    να ταλαντώνεται έντονα κατά τη διαδικασία selection.
    \item \textbf{F1-weighted} ζυγίζει κάθε κλάση με τη support της -
    είναι default για multi-class στο scikit-learn, και σταθερότερη
    μετρική για hyperparameter selection.
\end{itemize}
Το F1-macro κρατιέται ως \emph{μετρική αναφοράς} για να κρίνουμε αν το
μοντέλο απλώς «μαθαίνει τις μεγάλες κλάσεις» ή αν πιάνει και τις rare
attacks.

\section{Logistic Regression}\label{sec:q2_lr}

\textbf{Διαίσθηση.} Γραμμικός classifier, softmax για multi-class:
μία γραμμική λειτουργία ανά κλάση, normalized να αθροίζουν σε 1. Η
απόφαση χωρίζεται από \emph{υπερεπίπεδα} στο feature space. Καλός
baseline: ταχύτατο, ερμηνεύσιμο (coefficients), δίνει calibrated
probabilities. Περιορισμοί: υποθέτει γραμμική σχέση features
\(\to\) log-odds, ευαίσθητο σε scale, δεν πιάνει αλληλεπιδράσεις
features.

\textbf{Grid.} 16 combinations:
\[
\text{C} \in \{0.01,\,0.1,\,1,\,10\},\quad
\text{solver} \in \{\text{lbfgs},\,\text{saga}\},\quad
\text{max\_iter} \in \{200,\,500\}.
\]

\textbf{Αποτέλεσμα.} Και στα δύο σενάρια επιλέγονται τα \emph{ίδια}
hyperparameters: \(\text{C}=10\), \texttt{solver}=lbfgs,
\(\text{max\_iter}=500\). Η σύγκλιση των δύο σεναρίων δείχνει ότι το
validation set των 26.838 rows είναι αρκετά μεγάλο ώστε να συμφωνήσει
με τον CV μέσο όρο. Το metrics όπως το βλέπουμε στο test:

\begin{table}[H]
    \centering
    \begin{tabular}{lrrrrr}
        \toprule
        \textbf{Σενάριο} & \textbf{Accuracy} & \textbf{Precision\(_w\)}
        & \textbf{Recall\(_w\)} & \textbf{F1\(_w\)} & \textbf{F1\(_m\)} \\
        \midrule
        A (split) & 0,9595 & 0,9750 & 0,9595 & 0,9647 & 0,7810 \\
        B (CV)    & 0,9595 & 0,9750 & 0,9595 & 0,9647 & 0,7810 \\
        \bottomrule
    \end{tabular}
    \caption{Logistic Regression - test metrics.}\label{tab:q2_lr}
\end{table}

Το \(\text{F1-weighted}=0{,}965\) ακούγεται εξαιρετικό αλλά το
\(\text{F1-macro}=0{,}781\) αποκαλύπτει την αδυναμία: το LR
αποτυγχάνει στις non-linearly separable κλάσεις. Όπως αναμέναμε,
γραμμικά υπερεπίπεδα δεν αρκούν για να ξεχωρίσουν Web-Attack υπογραφές
μέσα σε BENIGN HTTP traffic.

\textbf{Συντελεστές του LR.} Το multinomial LR παράγει έναν πίνακα
συντελεστών shape \((n_\text{classes}, n_\text{features})\). Κάθε row
δείχνει πόσο σπρώχνει κάθε feature ένα sample \emph{προς} (θετικό
πρόσημο) ή \emph{μακριά από} (αρνητικό) την κλάση. Σε αντίθεση με τα
tree-based feature importances, οι LR συντελεστές δίνουν \emph{direction},
όχι μόνο magnitude - αυτό είναι το ερμηνεύσιμο πλεονέκτημα του LR.

\begin{figure}[H]
    \centering
    \includegraphics[width=0.95\textwidth]{q2_lr_coefficients.png}
    \caption{Logistic Regression (Σενάριο B) - heatmap των συντελεστών
    για τα top-15 features (κατά μέσο \(|\text{coef}|\) σε όλες τις
    κλάσεις). Έντονο κόκκινο = το feature οδηγεί strongly προς αυτή
    την κλάση· έντονο μπλε = αντενδείκνυται. Παράδειγμα: η
    \texttt{PortScan} έχει σαφώς θετικούς συντελεστές στα
    \texttt{Bwd Packet Length Min} και \texttt{Bwd Packets/s}
    (μικρά backward packets σε υψηλό rate - υπογραφή scanning).}
    \label{fig:q2_lr_coefficients}
\end{figure}

\section{Decision Tree}\label{sec:q2_dt}

\textbf{Διαίσθηση.} Ακολουθία \texttt{if/else} ερωτήσεων· σε κάθε
κόμβο επιλέγουμε feature+threshold που μεγιστοποιεί το
\emph{information gain} (μείωση impurity). Τα criteria που
συγκρίνουμε:
\begin{itemize}
    \item \textbf{Gini impurity}: \(1 - \sum_c p_c^2\). Πιθανότητα
    λανθασμένης ταξινόμησης αν βάλουμε τυχαία row στο πιο συχνό class.
    \item \textbf{Entropy}: \(-\sum_c p_c \log_2 p_c\). Information-
    theoretic measure - πόσα bits χρειάζονται για να περιγράψουμε την
    κλάση.
\end{itemize}
Regularization μέσω \texttt{max\_depth} (upper bound) και
\texttt{min\_samples\_split} (lower bound μεγέθους κόμβου). Δεν
απαιτείται scaling.

\textbf{Grid.} 24 combinations:
\[
\text{max\_depth}\in\{5,10,20,\texttt{None}\},\
\text{min\_samples\_split}\in\{2,10,50\},\
\text{criterion}\in\{\text{gini},\text{entropy}\}.
\]

\textbf{Αποτέλεσμα.} Σενάριο A επιλέγει
\(\text{max\_depth}=20\), \(\text{min\_samples\_split}=10\), entropy.
Σενάριο B επιλέγει \(\text{max\_depth}=\texttt{None}\) (unbounded),
\(\text{min\_samples\_split}=2\), entropy. Η διαφορά είναι μικρή σε
μετρικές απόδοσης:

\begin{table}[H]
    \centering
    \begin{tabular}{lrrrrr}
        \toprule
        \textbf{Σενάριο} & \textbf{Accuracy} & \textbf{Precision\(_w\)}
        & \textbf{Recall\(_w\)} & \textbf{F1\(_w\)} & \textbf{F1\(_m\)} \\
        \midrule
        A (split) & 0,9918 & 0,9930 & 0,9918 & 0,9920 & 0,9093 \\
        B (CV)    & 0,9921 & 0,9920 & 0,9921 & 0,9920 & 0,9097 \\
        \bottomrule
    \end{tabular}
    \caption{Decision Tree - test metrics.}\label{tab:q2_dt}
\end{table}

Το F1-macro βελτιώνεται κατά \(\sim 13\) ποσοστιαίες μονάδες σε σχέση
με το LR - ο DT πιάνει τις non-linear αλληλεπιδράσεις (π.χ. «μικρά
packets \emph{και} υψηλό SYN count») με multi-split paths.

\textbf{Feature importance.} Η \texttt{Destination Port} κυριαρχεί
συντριπτικά (\(\sim 0{,}39\) importance), ακολουθούμενη από
\texttt{Init\_Win\_bytes\_backward} (\(0{,}12\)), \texttt{Flow Duration}
(\(0{,}07\)), \texttt{Bwd Packet Length Max} (\(0{,}07\)) και
\texttt{min\_seg\_size\_forward} (\(0{,}07\)). Η κυριαρχία ενός
feature είναι χαρακτηριστική ενός μεμονωμένου δέντρου - το RF θα την
εξομαλύνει.

\textbf{Visualization του δέντρου.} Το πραγματικό best DT έχει
\(\text{max\_depth}=\texttt{None}\) (unbounded), οπότε η πλήρης
οπτικοποίηση είναι αδιάβαστη. Παρουσιάζουμε δύο εικόνες: (α) τα top-3
levels του actual best DT (Σχ.~\ref{fig:q2_dt_top3}), και (β) ένα
ξεχωριστό shallow tree \(\text{max\_depth}=3\) που εκπαιδεύσαμε μόνο
για διδακτική απεικόνιση (Σχ.~\ref{fig:q2_dt_shallow}). Σε κάθε
κόμβο εμφανίζονται: split rule (π.χ.
\texttt{Destination Port \(\le 80{,}5\)}), entropy, αριθμός samples
που κατέληξαν εκεί, κατανομή κλάσεων (\texttt{value=[\dots]}), και
majority class (color-coded).

\begin{figure}[H]
    \centering
    \includegraphics[width=0.95\textwidth]{q2_dt_visualization_top3.png}
    \caption{Decision Tree (Σενάριο B, actual best) - top 3 levels.
    Το πραγματικό δέντρο μπορεί να φτάνει σε πολύ μεγαλύτερο βάθος,
    αλλά τα πρώτα splits χωρίζουν τις περισσότερες κλάσεις.}
    \label{fig:q2_dt_top3}
\end{figure}

\begin{figure}[H]
    \centering
    \includegraphics[width=0.95\textwidth]{q2_dt_visualization_shallow.png}
    \caption{Standalone shallow Decision Tree (\(\text{max\_depth}=3\),
    entropy) εκπαιδευμένο ξεχωριστά για clean visualization. \emph{Όχι}
    το best model - απλώς δείχνει πώς θα μοιάζει ένα ολόκληρο δέντρο 3
    επιπέδων.}
    \label{fig:q2_dt_shallow}
\end{figure}

\section{Random Forest}\label{sec:q2_rf}

\textbf{Διαίσθηση.} Ensemble από πολλά δέντρα που διαφέρουν μεταξύ
τους μέσω:
\begin{enumerate}
    \item \textbf{Bagging}: κάθε δέντρο εκπαιδεύεται σε ένα bootstrap
    sample του training set (με replacement).
    \item \textbf{Feature randomness}: σε κάθε split, εξετάζεται μόνο
    ένα τυχαίο υποσύνολο features (default \(\sqrt{n}\)).
\end{enumerate}
Ένα μεμονωμένο δέντρο είναι high-variance· ο μέσος όρος πολλών
στατιστικά ανεξάρτητων δέντρων μειώνει τη variance χωρίς να αυξάνει
το bias. Επίσης, δεν υπερ-εμπιστεύεται κανένα μεμονωμένο feature.
Δεν απαιτείται scaling.

\textbf{Grid.} 18 combinations:
\[
\text{n\_estimators}\in\{50,100,200\},\
\text{max\_depth}\in\{10,20,\texttt{None}\},\
\text{min\_samples\_split}\in\{2,10\}.
\]

\textbf{Αποτέλεσμα.} Σενάριο A επιλέγει
\(\text{n\_estimators}=100\), \(\text{max\_depth}=20\),
\(\text{min\_samples\_split}=2\). Σενάριο B επιλέγει
\(\text{n\_estimators}=100\), \(\text{max\_depth}=\texttt{None}\),
\(\text{min\_samples\_split}=2\).

\begin{table}[H]
    \centering
    \begin{tabular}{lrrrrr}
        \toprule
        \textbf{Σενάριο} & \textbf{Accuracy} & \textbf{Precision\(_w\)}
        & \textbf{Recall\(_w\)} & \textbf{F1\(_w\)} & \textbf{F1\(_m\)} \\
        \midrule
        A (split) & 0,9920 & 0,9920 & 0,9920 & 0,9920 & 0,9116 \\
        B (CV)    & 0,9925 & 0,9920 & 0,9925 & 0,9921 & 0,9062 \\
        \bottomrule
    \end{tabular}
    \caption{Random Forest - test metrics.}\label{tab:q2_rf}
\end{table}

Το F1-weighted είναι οριακά καλύτερο σε B (0,9921 vs 0,9920), ενώ
το F1-macro είναι καλύτερο σε A (0,9116 vs 0,9062). Αυτό είναι
ενδιαφέρον: η CV επιλέγει με βάση F1-weighted (κατά convention στο
\texttt{GridSearchCV}), επομένως δεν εγγυάται τη βέλτιστη macro
απόδοση. Σε ένα production IDS που απαιτεί ισχυρή ανίχνευση rare
attacks, θα εναλλάσσαμε το scoring σε F1-macro.

\textbf{Feature importance.} Σε αντίθεση με τον DT, το RF κατανέμει
το importance: top-5 είναι \texttt{Destination Port} (\(0{,}073\)),
\texttt{Init\_Win\_bytes\_backward} (\(0{,}057\)),
\texttt{Fwd Packet Length Max} (\(0{,}045\)),
\texttt{Bwd Packets/s} (\(0{,}045\)),
\texttt{Packet Length Mean} (\(0{,}041\)). Δεν υπάρχει ένα feature-«μονάρχης»·
τα επόμενα 15 features έχουν παρόμοιες τιμές. Αυτή η εξισορρόπηση
είναι το signature του bagging/feature-randomness averaging.

\begin{figure}[H]
    \centering
    \includegraphics[width=0.9\textwidth]{q2_rf_feature_importance.png}
    \caption{Top-20 feature importances του Random Forest (Σενάριο B).}
    \label{fig:q2_rf_importance}
\end{figure}

\textbf{Visualization ενός voter του δάσους.} Το Σχ.~\ref{fig:q2_rf_one_tree}
δείχνει τα top-3 levels ενός μεμονωμένου δέντρου από το best RF
(δέντρο \#0 από τα 100). Σημειώνουμε ότι:
\begin{itemize}
    \item Το πρώτο split εδώ είναι \texttt{Bwd Packet Length Max} -
    \emph{διαφορετικό} από το πρώτο split του standalone DT
    (\texttt{Destination Port}). Αυτό είναι ακριβώς το diversity που
    επιδιώκουμε με bagging+feature-randomness.
    \item Αν δείχναμε ένα δεύτερο δέντρο (π.χ. \texttt{tree\_index=1}),
    η δομή θα ήταν ξανά διαφορετική.
    \item Η τελική πρόβλεψη του RF είναι \emph{majority vote} (ή average
    των per-tree probabilities) πάνω σε όλα τα 100 δέντρα. Άρα ένα
    μεμονωμένο tree δεν αντιπροσωπεύει την «απόφαση» του RF.
\end{itemize}

\begin{figure}[H]
    \centering
    \includegraphics[width=0.95\textwidth]{q2_rf_one_tree.png}
    \caption{Random Forest (Σενάριο B) - top 3 levels του πρώτου voter
    του δάσους (από τα 100 δέντρα συνολικά). Τα splits εδώ διαφέρουν
    από αυτά του standalone DT - σημείο της feature randomness που
    επιβάλλει το RF.}
    \label{fig:q2_rf_one_tree}
\end{figure}

\section{Συγκριτική Ανάλυση}\label{sec:q2_comparison}

\begin{table}[H]
    \centering
    \small
    \begin{tabular}{llrrrrr}
        \toprule
        \textbf{Model} & \textbf{Σενάριο} & \textbf{Acc} &
        \textbf{Prec\(_w\)} & \textbf{Rec\(_w\)} & \textbf{F1\(_w\)} &
        \textbf{F1\(_m\)} \\
        \midrule
        Logistic Regression & A & 0,9595 & 0,9750 & 0,9595 & 0,9647 & 0,7810 \\
        Logistic Regression & B & 0,9595 & 0,9750 & 0,9595 & 0,9647 & 0,7810 \\
        Decision Tree       & A & 0,9918 & 0,9930 & 0,9918 & 0,9920 & 0,9093 \\
        Decision Tree       & B & 0,9921 & 0,9920 & 0,9921 & 0,9920 & 0,9097 \\
        Random Forest       & A & 0,9920 & 0,9920 & 0,9920 & 0,9920 & \textbf{0,9116} \\
        Random Forest       & B & \textbf{0,9925} & 0,9920 &
        \textbf{0,9925} & \textbf{0,9921} & 0,9062 \\
        \bottomrule
    \end{tabular}
    \caption{Q2 - σύγκριση 6 παραλλαγών στο ίδιο test set
    (26.838 rows).}\label{tab:q2_comparison}
\end{table}

\begin{figure}[H]
    \centering
    \includegraphics[width=0.95\textwidth]{q2_comparison_f1.png}
    \caption{F1-weighted έναντι F1-macro και για τις 6 παραλλαγές.
    Η διαφορά F1\(_w\) - F1\(_m\) είναι η «υπογραφή του imbalance»:
    πόσο «χρωστά» το μοντέλο την απόδοσή του στις μεγάλες κλάσεις.}
    \label{fig:q2_comparison}
\end{figure}

\subsection{Απαντήσεις στις Τέσσερις Ερωτήσεις της Εκφώνησης}

\paragraph{1. Ποιος αλγόριθμος αποδίδει καλύτερα;}
Τα Decision Tree και Random Forest είναι σχεδόν ισόπαλα στο
F1-weighted (\(0{,}992\)) και στο F1-macro (\(\sim 0{,}91\)). Το
Random Forest κερδίζει οριακά στο accuracy (B: \(0{,}9925\)), όμως
\emph{το μεμονωμένο DT είναι καλύτερο στο F1-macro}. Η εξήγηση:
για αυτό το συγκεκριμένο dataset μετά το feature selection του Q1,
ένα καλά ρυθμισμένο δέντρο πιάνει όλα τα patterns· το RF averaging
δεν προσφέρει μεγάλη οριακή βελτίωση γιατί το πρόβλημα δεν είναι
high-variance-limited. Το LR υστερεί σημαντικά (F1-macro
\(=0{,}781\)) επειδή οι attack classes \emph{δεν} είναι γραμμικά
διαχωρίσιμες στο 47-feature space.

\paragraph{2. Δίνει το Grid Search με CV διαφορετικά hyperparameters σε σχέση με split-based;}
\textbf{LR}: \emph{ίδια} best params και στα δύο σενάρια -
το single validation set των \(\sim 27\)K rows είναι αρκετά μεγάλο
ώστε να συμφωνήσει με τη CV. \textbf{DT}: διαφορετικά
(A: max\_depth\(=20\), B: unbounded). \textbf{RF}: διαφορετικά
(A: max\_depth\(=20\), B: unbounded). Η CV τείνει να προτιμά
unbounded depth γιατί ο μέσος όρος 5 folds είναι σταθερότερος σε
ελαφρές διαφορές· ένα single val set μπορεί τυχαία να «τιμωρήσει» το
unbounded model αν κάποια rare-class rows πέσουν εκεί και overfittingιστούν.
\textbf{Γενικό μάθημα}: όταν το dataset είναι περιορισμένο ή οι
κλάσεις σπάνιες, η CV είναι πιο αξιόπιστη.

\paragraph{3. Ποιες attack types είναι οι πιο δύσκολες;}
Από τις row-normalized confusion matrices (Figs.~\ref{fig:q2_cm_rf},
\ref{fig:q2_cm_lr}):
\begin{itemize}
    \item \textbf{Web Attacks (Brute Force / XSS / SQL Injection)}:
    συγχέονται μεταξύ τους και με BENIGN. Αναπαράγουν normal HTTP
    traffic με μικρές υπογραφές.
    \item \textbf{Infiltration (36 rows)} και
    \textbf{Heartbleed (11 rows)}: λόγω περιορισμένων training data
    τα signatures είναι αναξιόπιστα. Συχνά κατηγοριοποιούνται ως
    άλλη high-frequency DoS κλάση.
    \item \textbf{Bot}: μπερδεύεται με BENIGN επειδή το command-and-
    control traffic μπορεί να μιμείται κανονικές HTTP/HTTPS calls.
\end{itemize}
Οι «εύκολες» κλάσεις είναι \texttt{DDoS, DoS Hulk, PortScan,
FTP-Patator, SSH-Patator} - έχουν πολύ διακριτές υπογραφές (rate
spikes, αποκλίνοντα packet sizes, καταιγίδες SYN χωρίς ACK).

\paragraph{4. Πώς επηρεάζει το class imbalance τα αποτελέσματα;}
Σε όλα τα μοντέλα ισχύει \(\text{F1}_w \gg \text{F1}_m\). Για το LR,
η διαφορά είναι \(0{,}965 - 0{,}781 = 0{,}184\) - τεράστια.
Για DT/RF μειώνεται σε \(\sim 0{,}08\), δηλώνοντας ότι τα non-linear
μοντέλα πιάνουν και τις rare classes (όχι τέλεια, αλλά πολύ καλύτερα).
Η επιλογή μετρικής εξαρτάται από την εφαρμογή:
\begin{itemize}
    \item IDS που θέλει να πιάνει \emph{κάθε} attack type \(\to\)
    F1-macro.
    \item Production filter που σταθμίζεται με τη συχνότητα traffic
    \(\to\) F1-weighted.
\end{itemize}

\begin{figure}[H]
    \centering
    \includegraphics[width=0.92\textwidth]{q2_rf_B_confusion.png}
    \caption{Random Forest (Σενάριο B) - row-normalized confusion
    matrix. Οι διαγώνιοι κοντά στο 1 δείχνουν per-class recall· οι
    off-diagonal entries δείχνουν confusion patterns.}
    \label{fig:q2_cm_rf}
\end{figure}

\begin{figure}[H]
    \centering
    \includegraphics[width=0.92\textwidth]{q2_lr_B_confusion.png}
    \caption{Logistic Regression (Σενάριο B) - row-normalized
    confusion matrix. Σημαντικά χειρότερη ποιότητα στις Web-Attack
    και Bot κλάσεις σε σχέση με το RF.}
    \label{fig:q2_cm_lr}
\end{figure}

\section{Σύνοψη Q2}\label{sec:q2_summary}

\begin{table}[H]
    \centering
    \begin{tabular}{lr}
        \toprule
        \textbf{Μέγεθος}                                     & \textbf{Τιμή} \\
        \midrule
        Subsample (cap \(=25\)K)                             & 134.187 rows  \\
        Train / Val / Test                                   & 80.511 / 26.838 / 26.838 \\
        Αλγόριθμοι                                           & LR, DT, RF    \\
        Σενάρια                                              & A (split), B (CV) \\
        LR grid combinations                                 & 16            \\
        DT grid combinations                                 & 24            \\
        RF grid combinations                                 & 18            \\
        Best model (F1-weighted)                             & RF B: 0,9921  \\
        Best model (F1-macro)                                & RF A: 0,9116  \\
        Best single-algorithm (F1-weighted, F1-macro combo)  & DT / RF παρόμοια \\
        Worst F1-macro                                       & LR: 0,7810    \\
        \bottomrule
    \end{tabular}
    \caption{Βασικά αποτελέσματα του Q2.}\label{tab:q2_summary}
\end{table}

%=============================================================================
% ΚΕΦΑΛΑΙΟ 5: ΕΡΩΤΗΜΑ 3 - CLUSTERING
%=============================================================================
\chapter{Ερώτημα 3: Συσταδοποίηση}

\textit{Το κεφάλαιο αυτό θα συμπληρωθεί μετά την υλοποίηση του Q3.
Αναμενόμενη δομή:}

\begin{itemize}
    \item \textbf{Preprocessing}: PCA για μείωση σε \(\sim 10\) components
    (κρίσιμο για distance-based αλγόριθμους).
    \item \textbf{Μοντέλα}: K-Means (με elbow/silhouette για επιλογή του
    \(K\)), Agglomerative Hierarchical, DBSCAN.
    \item \textbf{Αξιολόγηση}: inertia, silhouette score, Adjusted Rand
    Index έναντι των γνωστών labels (external validation).
    \item \textbf{Ερμηνεία}: κατά πόσο οι unsupervised συστάδες
    αντιστοιχούν στις γνωστές επιθέσεις.
\end{itemize}

%=============================================================================
% ΚΕΦΑΛΑΙΟ 6: ΣΥΜΠΕΡΑΣΜΑΤΑ
%=============================================================================
\chapter{Συμπεράσματα}

\textit{Το κεφάλαιο αυτό θα ολοκληρωθεί μετά τα Q2 και Q3. Έως τώρα, τα
βασικά συμπεράσματα του Q1 είναι:}

\begin{enumerate}
    \item \textbf{Το CIC-IDS-2017 απαιτεί προσεκτικό καθαρισμό} πριν
    οποιοδήποτε μοντέλο: 4.376 Infs, 255.236 exact duplicates και
    corrupted bytes στις ετικέτες των Web Attacks.

    \item \textbf{Η ανισορροπία είναι ακραία} (\(\sim 195\,000:1\)). Θα
    διαμορφώσει όλες τις επιλογές μοντελοποίησης στο Q2 (class\_weight,
    macro-F1, stratified CV) και αξιολόγησης στο Q3.

    \item \textbf{Το feature space είναι υπερ-παραμετροποιημένο}. Οι 78
    αρχικές αριθμητικές στήλες κρύβουν 22 near-duplicates και 8 constants.
    Μετά την αφαίρεσή τους μένουν 47 features με γνήσια διακριτική
    πληροφορία.

    \item \textbf{Η χρήση Parquet είναι κρίσιμη} για την αποδοτικότητα του
    αγωγού: Q2/Q3 φορτώνουν το καθαρισμένο σύνολο σε \(<2\) sec αντί για
    \(\sim 60\) sec.
\end{enumerate}

%=============================================================================
% ΠΑΡΑΡΤΗΜΑ
%=============================================================================
\appendix

\chapter{Οδηγίες Εγκατάστασης και Εκτέλεσης}

\section{Απαιτήσεις}

\begin{itemize}
    \item Python 3.10+ (δοκιμάστηκε σε 3.12).
    \item \(\sim\) 16 GB RAM (το full load απαιτεί \(\sim\) 4-6 GB
    μετά το \texttt{\_optimize\_memory}).
    \item \(\sim\) 3 GB ελεύθερου χώρου (raw CSV + Parquet cache).
\end{itemize}

\section{Εγκατάσταση}

\begin{tcolorbox}[myCustomStyle={Εγκατάσταση}]
\begin{lstlisting}[style=myBashStyle]
# clone και venv
git clone https://github.com/IBilba/network-intrusion-detection-ml.git
cd network-intrusion-detection-ml
python -m venv .venv
source .venv/bin/activate   # ή .venv\Scripts\activate στα Windows

# dependencies
pip install -r requirements.txt

# data (τοποθέτηση των 8 CSV στο data/)
# από Kaggle: https://www.kaggle.com/datasets/cicdataset/cicids2017
\end{lstlisting}
\end{tcolorbox}

\section{Εκτέλεση του Q1}

\begin{tcolorbox}[myCustomStyle={Run}]
\begin{lstlisting}[style=myBashStyle]
# Εκτέλεση end-to-end (γράφει figures, CSV, Parquet)
jupyter nbconvert --to notebook --execute notebooks/q1_eda.ipynb \
  --output q1_eda.ipynb
\end{lstlisting}
\end{tcolorbox}

Μετά την εκτέλεση, τα παραγόμενα αρχεία βρίσκονται σε:

\begin{itemize}
    \item \texttt{outputs/figures/}: \(5\) PNGs (label distribution,
    histograms, boxplots, heatmap).
    \item \texttt{outputs/results/}: \(7\) CSVs (column catalogue,
    describe, label distribution, correlation με label, top correlated
    pairs, removed features, final feature list).
    \item \texttt{data/cic\_ids\_2017\_clean.parquet}: το cached clean
    dataset για τα Q2/Q3.
\end{itemize}

\section{Αναπαραγωγιμότητα}

Όλες οι στοχαστικές λειτουργίες (sampling, shuffling, CV folds) θα
χρησιμοποιούν τον ίδιο \texttt{RANDOM\_STATE = 42} από το
\texttt{utils/helpers.py}. Με τον ίδιο seed και το ίδιο
σύνολο δεδομένων, όλοι οι αριθμοί του Πίνακα~\ref{tab:q1_summary} είναι
bit-exactly αναπαραγώγιμοι.

\end{document}
